% ==============================================================================
% ch:inovace — pět návrhů inovace sociálního dialogu v ČR
% ==============================================================================

Zhodnocení formulované v~kap.~\ref{sec:zhodnoceni} je~v~této kapitole transformováno do~pěti vymezených právních a~organizačních inovací sociálního dialogu v~\acloc{ČR}. Předložené návrhy nejsou koncipovány jako~uzavřený reformní balík, nýbrž jako~vzájemně komplementární intervence s~cílem zvýšit reálnou účinnost \acgen{KV} bez~výrazného nárůstu administrativní zátěže či celkového zdanění práce. Opatření se~opírají o~zahraniční zkušenosti doložené v~komparativních částech práce.~\cite{InovacniDoporuceni_ZahraniceKV}\par

\paragraph{Inovace 1: Aktivace legislativního rámce rozšiřování závaznosti odvětvových smluv}\label{par:inovace_ksvs}
Reaguje na~hluboký deficit v~pokrytí zaměstnanců (\SI{31}{\percent}--\SI{35}{\percent}) kolektivním vyjednáváním. Institut extenze závaznosti \acpgen{KSVS} dle~\mbox{§\,7} zákona o~\acloc{KV} (č.~2/1991~Sb.) v~české legislativě formálně existuje, avšak v~praxi je~téměř nevyužit. Vzorem pro~jeho oživení je rakouský model odvětvového vyjednávání s~vysokým pokrytím i~při nižší organizovanosti (kap.~\ref{par:at_model}). Navrhuje se:
\begin{itemize}[noitemsep]
  \item Standardizace datových podkladů pro~posuzování reprezentativnosti (přesné vymezení odvětvových hranic a~počtů členů).
  \item Zefektivnění notifikační procedury na~úrovni \acgen{MPSV} splňující předpoklad předběžného celkové formálního přezkumu.
  \item Zákonná povinnost pro~\acs{MPSV} publikovat podrobné, veřejně dostupné odůvodnění v~případě neudělení navrhované extenze.
\end{itemize}
Zefektivněním \mbox{§\,7} lze významně snížit transakční náklady na~straně navrhovatelů bez~nutnosti přistupovat k~politicky citlivým novelám samotného textu zákona. Opatření je~plně v~souladu se závazky směrnice o~adekvátních minimálních mzdách v~\acs{EU} a~doporučením \acs{ETUI} \emph{Road to~80~percent}.~\cite{eu_directive_2022_2041}~\cite{ETUI_RoadTo80_Muller_2025}\par
\newpage
\paragraph{Inovace 2: Odborové svazy jako spolutvůrci a~vykonavatelé aktivní politiky zaměstnanosti}\label{par:inovace_ghent}
Inovativní řešení řeší souběžně nízkou organizovanost odborů a~podfinancování \acgen{APZ}. Využívá inspiraci z~ghentského systému (kap.~\ref{par:dk_model}), kde se odborové organizace podílejí na rekvalifikačních programech, což posiluje jejich legitimitu a~členskou základnu. Navrhuje se:
\begin{itemize}[noitemsep]
  \item Umožnit reprezentativním odborovým centrálám (\ac{ČMKOS}, ASO) na základě smluvních partnerství s~Úřadem práce ČR garantovat a~organizovat specifické odvětvové rekvalifikační programy.
  \item Postupně navýšit finanční zdroje na \acacc{APZ} ze současných \SI{0{,}16}{\percent} na cílové \SI{1{,}0}{\percent}~\acs{HDP} v~horizontu pěti let, čímž dojde k~přiblížení se průměru států \aclgen{EU}.
  \item Saturovat financování této expanze zčásti ze sektorových bipartitních fondů (vzdělávání a~rozvoje) zakotvených přímo v~\acloc{KSVS}, po~vzoru osvědčených nizozemských fondů O\&O.
\end{itemize}
Tento krok poskytuje věcnou a~metodickou oporu kapacitám odborů v~oblasti celoživotního učení plynoucí z~poptávky firem po nové kvalifikaci v~důsledku automatizace a~digitalizace.\par

\paragraph{Inovace 3: Mzdová transparentnost nad rámec minimálního transpozičního standardu}\label{par:inovace_transparentnost}
Reaguje na transpozici směrnice EU 2023/970 o~posílení uplatňování zásady stejné odměny za stejnou práci. Doporučuje se překročit minimální transpoziční nároky ve~třech bodech:
\begin{itemize}[noitemsep]
  \item Zavést povinnost zveřejňování mzdových tarifů sjednaných v~\acloc{KSVS} v~rámci otevřeného rozhraní registru \acgen{MPSV}.
  \item Legislativně ukotvit jednotnou metodiku oceňování práce stejné hodnoty s~přímou referencí na klasifikační třídy v~rozšířených \acloc{KSVS}.
  \item Snížit práh pro povinné vykazování mzdových rozdílů mezi pohlavími na zaměstnavatele s~více než 50~zaměstnanci (oproti směrnicí doložené hranici 100~zaměstnanců).
\end{itemize}
Transparentní odměňování mírní mzdovou asymetrii a~podpoří agregátní poptávku.\par

\paragraph{Inovace 4: Integrace \acs{OSVČ} a~platformových pracovníků do rámce kolektivního vyjednávání}\label{par:inovace_svarc}
Adresuje rizika spojená se zastřeným zaměstnáváním, švarcsystémem a~prekérní povahou platformové práce. Navrhují se tři provázaná opatření:
\begin{itemize}[noitemsep]
  \item Navýšit analytické a~kontrolní kapacity \acgen{SÚIP} k~zefektivnění detekce zastřeného pracovního poměru za využití sektorových indikátorů ekonomické závislosti.
  \item Důsledná transpozice směrnice o~platformové práci s~vyvratitelnou právní domněnkou existence pracovního poměru v~souladu s~judikaturou \acgen{SDEU} (věc \emph{C-692/19 Yodel}).
  \item Explicitní legislativní rozšíření věcné a~personální působnosti \acpgen{KSVS} na samostatně činné osoby (solo-\acp{OSVČ}) vykazující prokazatelné znaky ekonomické závislosti, což odpovídá Pokynům Evropské komise z~roku \SI{2022}{\rok}.\par
\end{itemize}
Tento krok rozšiřuje vyjednávací bázi pro \ac{KV}, narovnává konkurenční prostředí na trhu a~zvyšuje odvodové příjmy sociálního zabezpečení v~souladu se zažitou koncepcí udržitelného rozvoje trhu práce.\par

\paragraph{Inovace 5: Digitální transparentnost a veřejný registr kolektivních smluv}\label{par:inovace_registr}
Odstraňuje informační asymetrie mezi sociálními partnery, statními institucemi a~výzkumnou obcí. Navrhuje se vybudovat veřejně přístupný, strojově čitelný registr všech odvětvových i~podnikových kolektivních smluv (vzor Belgie, Nizozemsko):
\begin{itemize}[noitemsep]
  \item Standardizovaný strukturovaný formát pro vykazování sjednaných mzdových tarifů, příplatků a~pracovní doby.
  \item Povinná elektronická notifikace nově uzavřených smluv s~funkcí automatické datové validace.
  \item Veřejně dostupné agregované statistické výstupy o~míře smluvního pokrytí dle klasifikace \acs{NACE} a~územního členění \acs{NUTS}.
\end{itemize}
Zřízení registru by minimalizovalo transakční náklady budoucích analýz a~poskytlo pevná data pro konstruktivní tripartitní vyjednávání. Technická proveditelnost je~vysoká, neboť \acgen{MPSV} již evidencí disponuje -- navrhované opatření by~bylo pouze uživatelskou a~funkční nadstavbou nad stávající infrastrukturou.\par

Závěrem lze konstatovat, že navržený soubor inovací kombinuje přímé věcné zásahy (1, 2, 4) s~podpůrnými transparentními nástroji (3, 5). Pouhé zavedení transparentnosti bez~mechanismů na zvýšení pokrytí kolektivními smlouvami by nepomohlo překonat institucionální slabost sociálního dialogu v~\acloc{ČR}. Naopak extenze bez~analytické opory a~přehledného digitálního registru by znesnadňovaly průběžnou kontrolu dopadů. Synergického a~konzistentního přínosu je proto možné dosáhnout pouze souběžnou a~provázanou realizací obou těchto rovin.\par
